\documentclass[10pt,a4paper]{amsart}
\usepackage[margin=19mm]{geometry}
\usepackage{amsmath,amssymb,mathtools}
\usepackage[scr=rsfs,cal=euler]{mathalfa}
\usepackage{xfrac}
\usepackage[unicode,hidelinks,bookmarks=false]{hyperref}
\newcommand{\ie}{i.e.\ }
\newcommand{\cf}{cf.\ }
\newcommand{\resp}{respectively\ }
\newcommand{\Subsection}{\subsection}
\providecommand{\nobreakdash}{\nobreak\hbox{-}}
\setlength{\emergencystretch}{2em}
\multlinegap0pt
\usepackage{amssymb}
\usepackage{mathtools}
\newtheorem{theorem}{Theorem}
\newtheorem{lemma}{Lemma}
\title{Complete Simplicial Fans, Stanley--Reisner Rings, and equivariant h-polynomial}
\author{Tao Gui}
\date{Source: arXiv:2605.15061v1 (CC BY 4.0)}
\begin{document}
\maketitle

\noindent\emph{Source and licence.} Complete Simplicial Fans, Stanley--Reisner Rings, and equivariant h-polynomial. Version arXiv:2605.15061v1 (CC BY 4.0), distributed by arXiv under the Creative Commons Attribution 4.0 International licence, http://creativecommons.org/licenses/by/4.0/, as recorded in the arXiv licence metadata for this exact version and shown on the abstract page https://arxiv.org/abs/2605.15061v1. Full licence notice: \texttt{licenses/arxiv-2605.15061-CC-BY-4.0.txt}.\par\smallskip

\noindent\textit{Editorial scope.} Counted unit: the article's theorem thm-main (source0.tex lines 113--119). The article's complete original proof of it is delivered with it (source0.tex lines 339--387, one \texttt{proof} environment of 49 lines, which the article places away from the statement). No mathematics is written, repaired or re-derived: the statement and the proof are the article's own lines, in the article's own notation, and no retained line is substituted or re-flowed.  Retained uncounted as the source-local context the proof needs: lines 192--194; lines 196--198; lines 241--243; lines 248--252; lines 260--262.  Nothing was omitted from any retained span, so the package carries no omission record.  No retained line contains a citation, so the wrapper carries no bibliography and no bibliography entry.  No retained line contains a figure environment, an image-inclusion command or drawing code, so no image is delivered and the package declares no asset dependency.  Editorial formatting only, in addition: an AMS \texttt{amsart} preamble that restates the article's own declarations and macros used by the retained lines, namely the environments \texttt{theorem}, \texttt{lemma} on separate counters, together with the standard packages those lines need.  The retained environments print in this document as Theorem 1, Lemma 1 in order of appearance, whereas the article prints the same items with the numbering of its own full document.  The complete native source of the article as distributed in the arXiv e-print for this exact version is bundled unchanged as \texttt{sources/200628/source0.tex}.
\begin{equation} \label{eq-decom of SR}
   \mathcal{A}(\Sigma)=\bigoplus_{\sigma \in \Sigma} \mathcal{A}(\Sigma)_\sigma, 
\end{equation}

\begin{equation} \label{eq-summand of SR}
\mathcal{A}(\Sigma)_\sigma=x_\sigma \cdot \mathbb{C}[x_{\rho_i} \mid \rho_i \in \sigma(1)]
\end{equation}

\begin{equation} \label{eq-Giso}
  \mathcal{A}(\Sigma)\cong \mathcal{A}\otimes_{\mathbb{R}}\overline{\mathcal{A}(\Sigma)}.  
\end{equation}

\begin{lemma} \label{lem-sym}
If $V$ is a $n$-dimensional vector space and $g \in \text{End}(V)$ is a linear transformation, then graded character of $g$ on the symmetric algebra $Sym(V)$ is given by 

$$\sum_{k=0}^{\infty} \text{Tr}(g|_{Sym^k(V)}) t^k = \frac{1}{\det(I - tg)}.$$
\end{lemma}

\begin{equation} \label{eq-PD}
 P_G(t, g)=t^d P_G(t^{-1}, g)   
\end{equation}

\begin{theorem} \label{thm-main}
For a finite group $G$ acting linearly on $V$ preserving a complete simplicial fan $\Sigma$, the equivariant Hilbert series of $\overline{\mathcal{A}(\Sigma)}$ is given by\footnote{Although the expression on the right hand side of \eqref{eq-general} contains fractions, an expression involving only integers can be obtained by allowing $\sigma$ to vary over a choice of representatives for the $G$-orbits of $\Sigma$.}
\begin{equation} \label{eq-general}
   \operatorname{R}_{\overline{\mathcal{A}(\Sigma)}}(t)=\sum_{\sigma \in \Sigma} \frac{\left|G_{\sigma }\right|}{\left|G\right|} \operatorname{Ind}_{G_{\sigma}}^{G}\left(\operatorname{det}(tI-\psi_\sigma)\right),
\end{equation}
where $\operatorname{det}(tI-\psi_\sigma)$ is viewed as a $\mathbb{C}[t]$-valued class function on $G_{\sigma }$ whose value is $\operatorname{det}(tI-\psi_\sigma(g))$ when evaluates at $g \in G$.
\end{theorem}

\begin{proof}[Proof of Theorem \ref{thm-main}]
Our strategy relies on three main steps: decomposing the Stanley--Reisner ring $\mathcal{A}(\Sigma)$ into $G$-invariant pieces, and then passing to the augmented quotient $\overline{\mathcal{A}(\Sigma)}$ using \eqref{eq-Giso}, 
and applying equivariant Poincaré duality (the algebraic manifestation of the Dehn-Sommerville equations) to align the polynomial degrees.

We want to evaluate the graded character (trace) of $g \in G$. Since $G$ permutes the cones in $\Sigma$, $g$ maps the summand $\mathcal{A}(\Sigma)_\sigma$ in \eqref{eq-decom of SR} to $\mathcal{A}(\Sigma)_{g\sigma}$. Thus, the only non-zero trace contributions come from the set of cones fixed setwisely by $g$, denoted $\Sigma^g$. 

For a invariant cone $\sigma \in \Sigma^g$, $g$ permutes its extreme rays. Because $\Sigma$ is a simplicial fan, the rays of $\sigma$ are linearly independent and form a basis for $V_\sigma$. Therefore, the permutation representation of $g$ on the variables $\{x_i \mid \rho_i \in \sigma(1)\}$ is isomorphic to the linear action of $g$ on $V_\sigma$. 

The generator $x_\sigma$ in \eqref{eq-summand of SR} is fixed by $g$ (since permutations leave the product of all variables invariant) and contributes $t^{\dim \sigma}$. The symmetric algebra $\mathbb{C}[x_{\rho_i} \mid \rho_i \in \sigma(1)]$ in \eqref{eq-summand of SR} contributes $1 / \operatorname{det}(I - t g_{V_\sigma})$ by Lemma \ref{lem-sym}. Summing over all invariant cones yields the graded character of $g \in G$ for the full Stanley-Reisner ring:
$$\operatorname{tr}(g \mid \mathcal{A}(\Sigma), t) = \sum_{\sigma \in \Sigma^g} \frac{t^{\dim \sigma}}{\operatorname{det}(I - t g_{V_\sigma})}.$$

Consider the isomorphism of graded $G$-modules in \eqref{eq-Giso}:
$$\mathcal{A}(\Sigma)\cong \mathcal{A}\otimes_{\mathbb{R}}\overline{\mathcal{A}(\Sigma)}.$$
Since the base field is $\mathbb{R}$, by Lemma \ref{lem-sym}, the graded character of $\mathcal{A}\cong\operatorname{Sym}(V^*)$ is $1 / \operatorname{det}(I - t g_{V^*}) = 1 / \operatorname{det}(I - t g_V)$  . Taking graded characters on both sides, we get the graded character of the Artinian reduction:
\begin{equation} \label{eq-cha from Giso}
  P_G(t, g) = \operatorname{det}(I - t g_V) \sum_{\sigma \in \Sigma^g} \frac{t^{\dim \sigma}}{\operatorname{det}(I - t g_{V_\sigma})}. 
\end{equation}

By Maschke's theorem, $V \cong V_\sigma \oplus (V/V_\sigma)$ as $g$-modules, which implies $\operatorname{det}(I - t g_V) = \operatorname{det}(I - t g_{V_\sigma}) \operatorname{det}(I - t g_{V/V_\sigma})$. This simplifies \eqref{eq-cha from Giso} to
\begin{equation} \label{eq-cha over inva cones}
P_G(t, g) = \sum_{\sigma \in \Sigma^g} t^{\dim \sigma} \operatorname{det}(I - t g_{V/V_\sigma}).
\end{equation} 

For an operator $A$ of dimension $k$, standard linear algebra dictates that $t^k \operatorname{det}(I - t^{-1}A) = \operatorname{det}(tI - A)$. Let's apply the palindromic property \eqref{eq-PD} to our sum \eqref{eq-cha over inva cones}:
\begin{equation} \label{eq-reverse}
    \begin{aligned}
P_G(t, g)=t^d P_G(t^{-1}, g) &= t^d \sum_{\sigma \in \Sigma^g} t^{-\dim \sigma} \operatorname{det}(I - t^{-1} g_{V/V_\sigma}) \\
&= \sum_{\sigma \in \Sigma^g} t^{d - \dim \sigma} \operatorname{det}(I - t^{-1} g_{V/V_\sigma}) \\
&= \sum_{\sigma \in \Sigma^g} \operatorname{det}(t I - g_{V/V_\sigma}),
\end{aligned} 
\end{equation}
where the last equality holds because $g_{V/V_\sigma}$ operates on a space of dimension $d - \dim(\sigma)$. 

We have established the character evaluated at a specific $g \in G$ as a sum over its fixed cones. The target formula \eqref{eq-general} is expressed globally using induced representations. Let's evaluate the target formula at $g$:
\begin{equation} \label{eq-eva at g}
\operatorname{R}_{\overline{\mathcal{A}(\Sigma)}}(t)(g) = \sum_{\sigma \in \Sigma} \frac{|G_\sigma|}{|G|} \operatorname{Ind}_{G_\sigma}^G \left(\operatorname{det}(tI-\psi_\sigma)\right)(g).    
\end{equation}  

Recall the Frobenius's formula for a general induced character: $$\operatorname{Ind}_{H}^G(\chi)(g) = \frac{1}{|H|} \sum_{h \in G, \ h^{-1}gh \in H} \chi(h^{-1}gh),$$
where $\chi$ is any character of a subgroup $H\subset G$.
Substituting this into the sum \eqref{eq-eva at g}, we have
\begin{equation} \label{eq-double}
\operatorname{R}_{\overline{\mathcal{A}(\Sigma)}}(t)(g) = \frac{1}{|G|} \sum_{\sigma \in \Sigma} \sum_{\substack{h \in G \\ h^{-1}gh \in G_\sigma}} \operatorname{det}(tI - \psi_\sigma(h^{-1}gh)_{V/V_\sigma})   
\end{equation}

Let $\tau = h\sigma$. The condition $h^{-1}gh \in G_\sigma$ means $gh\sigma = h\sigma$, which is exactly $g\tau = \tau$, meaning $\tau \in \Sigma^g$. Furthermore, the conjugation by $h$ intertwines the representation on $V/V_\sigma$ with $V/V_\tau$, so $\operatorname{det}(tI - \psi_\sigma(h^{-1}gh)_{V/V_\sigma}) = \operatorname{det}(tI - \psi_\sigma(g)_{V/V_\tau})$. Hence the double sum in \eqref{eq-double} re-indexes over all $h \in G$ and all $\tau \in \Sigma^g$:
$$\operatorname{R}_{\overline{\mathcal{A}(\Sigma)}}(t)(g) = \frac{1}{|G|} \sum_{h \in G} \sum_{\tau \in \Sigma^g} \operatorname{det}(tI - g_{V/V_\tau}) = \sum_{\tau \in \Sigma^g} \operatorname{det}(tI - g_{V/V_\tau}).$$
This precisely matches \eqref{eq-reverse}, which complete the proof of Theorem \ref{thm-main}.
\end{proof}

\end{document}
